% Figure: thermal-runaway heat-balance diagram. Heat-generation curve rises
% exponentially (Arrhenius) with temperature; heat-removal line is linear in
% (T - T_coolant). Intersections are operating points: a stable low point, an
% unstable middle point (tangency / ignition), and the runaway branch above.
% Two removal lines shown for different cooling capacities.
\begin{figure}[htbp]
\centering
\begin{tikzpicture}
\begin{axis}[
  width=12cm, height=7.5cm,
  xlabel={temperature $T$},
  ylabel={heat rate $\dot{q}$},
  xmin=0, xmax=10, ymin=0, ymax=10,
  axis lines=left,
  legend pos=north west,
  legend cell align=left,
  tick label style={font=\scriptsize},
  label style={font=\small},
  legend style={font=\scriptsize},
  xtick=\empty, ytick=\empty,
  samples=120, domain=0:10,
  clip=false,
]
% generation: Arrhenius-like sigmoid (heat release rises then saturates as
% reactant depletes). Use a logistic-shaped curve so it stays bounded.
\addplot[engred, very thick] {9.2/(1+exp(-1.5*(x-5.2)))};
\addlegendentry{generation $\dot{q}_g(T)$}
% removal line 1: adequate cooling, steep slope -> single low stable point
\addplot[engblue, very thick] {1.15*(x-1.0)};
\addlegendentry{removal $\dot{q}_r(T)$, high cooling}
% removal line 2: marginal cooling, shallow slope -> three intersections
\addplot[engamber, very thick, dashed] {0.62*(x-0.8)};
\addlegendentry{removal $\dot{q}_r(T)$, low cooling}
% mark intersections on the marginal-cooling case (approximate, illustrative)
\fill[engamber] (axis cs:1.55,0.46) circle (2.0pt);
\node[font=\scriptsize, engamber, anchor=north west] at (axis cs:1.55,0.5) {stable};
\fill[engamber] (axis cs:4.5,2.29) circle (2.0pt);
\node[font=\scriptsize, engamber, anchor=south east] at (axis cs:4.5,2.29) {unstable (ignition)};
\fill[engamber] (axis cs:8.4,4.71) circle (2.0pt);
\node[font=\scriptsize, engamber, anchor=south east] at (axis cs:8.4,4.71) {runaway branch};
\end{axis}
\end{tikzpicture}
\caption[Thermal-runaway heat-balance diagram]{The heat-balance diagram
that diagnoses thermal runaway. The generation curve $\dot{q}_g(T)$ rises
steeply with temperature (Arrhenius kinetics in a reactor, leakage
current in a chip, side reactions in a cell) and saturates as the driving
reactant depletes. The removal rate $\dot{q}_r(T)$ is close to linear in
the temperature above coolant. Where the two cross is a steady operating
point: it is stable if the removal curve cuts the generation curve from
below ($d\dot{q}_r/dT > d\dot{q}_g/dT$) and unstable otherwise. With ample
cooling the steep removal line meets the generation curve at a single low
stable point. With marginal cooling the shallow line produces a low
stable point, an unstable ignition point, and a high runaway branch; a
disturbance past the ignition point drives the system up to the runaway
branch. The design rule is to keep the removal line steep enough that no
unstable intersection sits inside the operating envelope.}
\label{fig:vol04ch05:runaway-balance}
\end{figure}
